% Preâmbulo comum aos documentos do Grupo VI (classe abnTeX2, normas ABNT).
% Compilar com Tectonic (XeTeX): ver gerar_pdf.sh.
\documentclass[12pt,openright,oneside,a4paper,chapter=TITLE,english,brazil]{abntex2}

% ---------------------------------------------------------------- Fontes
\usepackage{fontspec}
\setmainfont{texgyretermes}[
  Extension=.otf,
  UprightFont=*-regular, BoldFont=*-bold,
  ItalicFont=*-italic, BoldItalicFont=*-bolditalic]
\setmonofont{DejaVuSansMono}[
  Extension=.ttf, BoldFont=*-Bold, Scale=0.82]
\setsansfont{texgyreheros}[
  Extension=.otf,
  UprightFont=*-regular, BoldFont=*-bold,
  ItalicFont=*-italic, BoldItalicFont=*-bolditalic]
\usepackage{amsmath}

% ---------------------------------------------------------------- Pacotes
\usepackage{indentfirst}
\usepackage{graphicx}
\usepackage{xcolor}
\usepackage{booktabs}
\usepackage{array}
\usepackage{tabularx}
\usepackage{longtable}
\usepackage{xltabular}
\usepackage{multirow}
\usepackage[section]{placeins}
\usepackage{enumitem}
\usepackage{listings}
\usepackage{tikz}
\usepackage{adjustbox}
\usetikzlibrary{arrows.meta,positioning,fit,calc,shapes.geometric,backgrounds}
\usepackage[alf,abnt-emphasize=bf]{abntex2cite}

\graphicspath{{../../backend/resultados/}}

% ---------------------------------------------------------------- Aparência
\definecolor{cinzacodigo}{gray}{0.96}
\definecolor{cinzaborda}{gray}{0.75}
\lstset{
  language=Python,
  basicstyle=\ttfamily\small,
  keywordstyle=\bfseries,
  commentstyle=\itshape\color{black!65},
  stringstyle=\color{black!80},
  backgroundcolor=\color{cinzacodigo},
  frame=single, rulecolor=\color{cinzaborda},
  framesep=4pt, xleftmargin=6pt, xrightmargin=6pt,
  columns=fullflexible, keepspaces=true,
  showstringspaces=false, breaklines=true,
  numbers=none, aboveskip=0.8em, belowskip=0.8em,
}
\lstdefinestyle{json}{language={}, basicstyle=\ttfamily\footnotesize, keywordstyle={}}
\renewcommand{\lstlistingname}{Código}
\AtBeginDocument{\renewcommand{\thelstlisting}{\arabic{lstlisting}}}
\renewcommand{\lstlistlistingname}{Lista de códigos}

\setlist{itemsep=0.2em, topsep=0.3em}
\setlist[itemize]{label=--}
\newcolumntype{L}[1]{>{\raggedright\arraybackslash}p{#1}}
\newcolumntype{Y}{>{\raggedright\arraybackslash}X}
\newcolumntype{C}[1]{>{\centering\arraybackslash}p{#1}}
\renewcommand{\arraystretch}{1.25}
\setlength{\arrayrulewidth}{0.5pt}
% Quadros (ABNT): texto, bordas fechadas. Tabelas (IBGE): dados numéricos, laterais abertas.
% Definição do ambiente quadro conforme a documentação do abnTeX2.
\newcommand{\quadroname}{Quadro}
\newcommand{\listofquadrosname}{Lista de quadros}
\newcommand{\quadroautorefname}{Quadro}
\newfloat[chapter]{quadro}{loq}{\quadroname}
\newlistof{listofquadros}{loq}{\listofquadrosname}
\newlistentry{quadro}{loq}{0}
\setfloatadjustment{quadro}{\centering}
\counterwithout{quadro}{chapter}
\renewcommand{\cftquadroname}{\quadroname\space}
\renewcommand*{\cftquadroaftersnum}{\hfill--\hfill}
\setfloatlocations{quadro}{hbtp}

% Fonte de figuras e tabelas produzidas pelo grupo
\newcommand{\fonteautores}{\fonte{Elaborado pelos autores (2026).}}
% Diagrama TikZ: \diagrama{arquivo} reduz a figura à largura do texto, se preciso
\newcommand{\diagrama}[1]{\begin{adjustbox}{max width=\textwidth}\input{figuras/#1}\end{adjustbox}\par}
% Nome de arquivo, função ou campo do código
% (permite quebra de linha depois de \_ e de /)
\newcommand{\codsublinhado}{\textunderscore\allowbreak}
\newcommand{\cod}[1]{{\ttfamily\let\_\codsublinhado#1}}

% ---------------------------------------------------------------- PDF
\makeatletter
\hypersetup{
  pdftitle={\@title},
  pdfauthor={Grupo VI -- Sistemas Distribuídos -- UNIFESSPA},
  pdfsubject={Agregador de Preços de Eletrônicos},
  pdfcreator={LaTeX com abnTeX2},
  colorlinks=false, hidelinks, bookmarksdepth=4}
\makeatother

\setlength{\parindent}{1.25cm}
\setlength{\parskip}{0.2cm}
\setlength{\emergencystretch}{3em}

% Títulos em Times, negrito (seções primárias em caixa alta, como pede a NBR 6024)
\renewcommand{\ABNTEXchapterfont}{\rmfamily\bfseries}
\renewcommand{\ABNTEXchapterfontsize}{\normalsize}
\renewcommand{\ABNTEXsectionfont}{\rmfamily\bfseries}
\renewcommand{\ABNTEXsectionfontsize}{\normalsize}
\renewcommand{\ABNTEXsubsectionfont}{\rmfamily\bfseries}
\renewcommand{\ABNTEXsubsectionfontsize}{\normalsize}
\renewcommand{\ABNTEXsubsubsectionfont}{\rmfamily\itshape}
\renewcommand{\ABNTEXsubsubsectionfontsize}{\normalsize}

% ---------------------------------------------------------------- Capa e folha de rosto
% Versões próprias (NBR 14724): instituição, autores, título, subtítulo, local e ano.
\newcommand{\subtitulodoc}{}
\newcommand{\definirsubtitulo}[1]{\renewcommand{\subtitulodoc}{#1}}
\makeatletter
\newcommand{\capagrupo}{%
  \begin{capa}
    \centering
    {\bfseries\imprimirinstituicao\par}
    \vfill
    {\bfseries\large\MakeUppercase{\imprimirtitulo}\par}
    \ifx\subtitulodoc\@empty\else
      \vspace{0.4cm}{\large\subtitulodoc\par}\fi
    \vfill
    {\imprimirlocal\par}
    {\imprimirdata\par}
    \vspace*{0.5cm}
  \end{capa}}
\newcommand{\folhaderostogrupo}{%
  \begin{folhaderosto}
    \centering
    \vspace*{\fill}
    {\bfseries\large\MakeUppercase{\imprimirtitulo}\par}
    \ifx\subtitulodoc\@empty\else
      \vspace{0.4cm}{\large\subtitulodoc\par}\fi
    \vspace*{\fill}
    \hspace{.45\textwidth}%
    \begin{minipage}{.5\textwidth}
      \SingleSpacing\small
      \imprimirpreambulo
      \par\vspace{0.6cm}
      % nomes na mesma linha, separados por vírgula (\\ de \autor vira ",")
      \textbf{Discentes:} {\def\\{,}\imprimirautor}.\par
      \vspace{0.6cm}
      \textbf{\imprimirorientadorRotulo}~\imprimirorientador
    \end{minipage}\par
    \vspace*{\fill}
    {\imprimirlocal\par}
    {\imprimirdata\par}
    \vspace*{0.5cm}
  \end{folhaderosto}}
\makeatother
\newcommand{\listafiguras}{%
  \pdfbookmark[0]{\listfigurename}{lof}%
  \listoffigures*%
  \cleardoublepage}
\newcommand{\listatabelas}{%
  \pdfbookmark[0]{\listtablename}{lot}%
  \listoftables*%
  \cleardoublepage}
% Página pré-textual com título centralizado e sem número de seção
\newcommand{\paginapretextual}[1]{%
  \pdfbookmark[0]{#1}{\detokenize{#1}}%
  \begin{center}{\ABNTEXchapterfont\MakeUppercase{#1}}\end{center}\vspace{0.8cm}}
\newcommand{\listaquadros}{%
  \pdfbookmark[0]{\listofquadrosname}{loq}%
  \listofquadros*%
  \cleardoublepage}
% Legenda de quadro fora de float (para quadros longos com xltabular)
\newfixedcaption{\legendaquadro}{quadro}
\newcommand{\sumario}{%
  \pdfbookmark[0]{\contentsname}{toc}%
  \tableofcontents*%
  \cleardoublepage}

% Dados da capa e da folha de rosto (comuns aos três documentos).
% Autores na ordem do enunciado (Grupo VI). Para completar os nomes, edite apenas \autor abaixo.
\instituicao{%
  UNIVERSIDADE FEDERAL DO SUL E SUDESTE DO PARÁ -- UNIFESSPA
  \par
  FACULDADE DE SISTEMAS DE INFORMAÇÃO -- FACSI
  \par
  CURSO DE SISTEMAS DE INFORMAÇÃO}
\autor{%
  Youri\\
  José Miguel\\
  Mateus Leitão\\
  Michelangelo\\
  Ângelo}
\orientador[Professor:]{Gleison de Oliveira Medeiros}
\local{Marabá -- PA}
\data{2026}
\newcommand{\naturezatrabalho}{Trabalho~1 da disciplina Sistemas Distribuídos -- Engenharia de Requisitos e Desenvolvimento de um Agregador Concorrente (Tema~6: Agregador de Preços de Eletrônicos), apresentado ao Curso de Sistemas de Informação da Faculdade de Sistemas de Informação da Universidade Federal do Sul e Sudeste do Pará. Grupo~VI.}

\titulo{Agregador de Preços de Eletrônicos}
\definirsubtitulo{Documento de Requisitos}
\preambulo{\naturezatrabalho}

\begin{document}
\selectlanguage{brazil}
\frenchspacing

% ================================================================= PRÉ-TEXTUAIS
\pretextual
\capagrupo
\folhaderostogrupo

\paginapretextual{Histórico de revisões}
\begin{center}\small
  \begin{tabularx}{\textwidth}{@{}l l Y@{}}
    \toprule
    \textbf{Versão} & \textbf{Data} & \textbf{Alterações} \\
    \midrule
    1.0 & nov. 2026 & Versão inicial: problema, objetivo, escopo, atores, requisitos, restrições, casos de uso UC01--UC04 e critérios de aceitação CA01--CA15. \\
    1.1 & nov. 2026 & Interface web (gateway HTTP como cliente TCP do agregador), caso de uso UC05 e critérios CA16--CA18. \\
    1.2 & nov. 2026 & Roteiro da demonstração organizado pelos cinco experimentos do enunciado (UC04, CA15). \\
    1.3 & nov. 2026 & Requisitos reduzidos aos oito RF e cinco RNF que seguem os exemplos das seções 5.2 e 5.3 do enunciado. \\
    1.4 & 15 nov. 2026 & Edição em \LaTeX{} segundo as normas da ABNT. \\
    \bottomrule
  \end{tabularx}
\end{center}
\cleardoublepage

\listafiguras
\listaquadros
\sumario

% ================================================================= TEXTUAIS
\textual

\chapter{Introdução}

Este documento especifica os requisitos do \textbf{Agregador de Preços de Eletrônicos}, desenvolvido pelo Grupo~VI como Trabalho~1 da disciplina Sistemas Distribuídos (Tema~6). A organização segue a seção~5 do enunciado \cite{enunciado2026}: análise do problema (descrição, objetivo e escopo), atores, requisitos funcionais e não funcionais, restrições, casos de uso e critérios de aceitação. A forma de apresentar os requisitos -- identificador, enunciado no formato ``o sistema deverá'', prioridade e rastreabilidade até o código -- segue \citeonline{sommerville2015}.

Os diagramas de casos de uso, de arquitetura e de sequência estão no documento de \textbf{Modelagem}; a implementação, os experimentos e os resultados estão no \textbf{Relatório Técnico}. O protocolo de aplicação está especificado em \cod{docs/03\_protocolo.md} e o plano de testes em \cod{docs/05\_plano\_verificacao\_testes.md}.

% -----------------------------------------------------------------
\chapter{Análise do problema}

\section{Descrição do problema}

Quem compra eletrônicos (placas de vídeo, notebooks, monitores, SSDs, smartphones) costuma consultar várias lojas para achar o melhor preço. Fazer isso manualmente é lento. Além disso, o preço anunciado nem sempre é o custo final, porque frete, garantia e reputação do produto variam de loja para loja.

Plataformas de comparação resolvem isso consultando várias fontes independentes. Cada consulta passa pela rede e tem uma latência própria. Se as consultas forem feitas \textbf{uma após a outra}, o tempo de resposta cresce com o número de lojas; se forem feitas \textbf{ao mesmo tempo}, o tempo total tende ao da loja mais lenta \cite{tanenbaum2007}.

\section{Objetivo do sistema}

Desenvolver uma aplicação cliente-servidor, em Python, com \textbf{sockets TCP} e \textbf{multithreading}, que:

\begin{itemize}
  \item receba do usuário a pesquisa de um eletrônico;
  \item consulte pelo menos quatro lojas independentes, cada uma um processo servidor;
  \item consolide, ordene e destaque a melhor oferta;
  \item permita comparar experimentalmente a estratégia \textbf{sequencial} com a \textbf{concorrente}.
\end{itemize}

\section{Escopo}

O \autoref{tab:escopo} delimita o que faz e o que não faz parte do sistema.

\begin{quadro}[!htbp]
  \centering\small
  \caption{Escopo do sistema}\label{tab:escopo}
  \begin{tabularx}{\textwidth}{|Y|Y|}
    \hline
    \textbf{Dentro do escopo} & \textbf{Fora do escopo} \\ \hline
    Cliente de linha de comando (interativo e por argumentos). & Integração com lojas reais, preços reais ou APIs comerciais. \\ \hline
    Interface web (Vite + React, em \cod{frontend/}) que usa um gateway HTTP local (\cod{backend/gateway\_http.py}) como cliente TCP do agregador. & Autenticação, carrinho de compras e pagamento. \\ \hline
    Agregador TCP que atende vários clientes ao mesmo tempo. & Persistência em banco de dados, cache e replicação. \\ \hline
    Quatro serviços de loja com catálogos simulados (arquivos JSON). & Comunicação entre máquinas diferentes (suportada pela configuração, mas não exigida nem testada). \\ \hline
    Consultas sequenciais e concorrentes (threads). & \\ \hline
    Ordenação por preço total e por custo-benefício (evolução do Tema~6). & \\ \hline
    Tratamento de falhas: loja fora do ar, queda durante a consulta, timeout e mensagens inválidas. & \\ \hline
    Registro do tempo de cada pesquisa e scripts para os cinco experimentos, com tabelas e gráficos. & \\ \hline  \end{tabularx}
  \fonteautores
\end{quadro}

\section{Funcionalidades mínimas do Tema 6}

O enunciado define, para o Tema~6, as funcionalidades mínimas e uma possibilidade de evolução \cite{enunciado2026}. O \autoref{qua:tema6} mostra o requisito funcional que atende cada uma.

\begin{quadro}[!htbp]
  \centering\small
  \caption{Funcionalidades do Tema 6 e requisitos correspondentes}\label{qua:tema6}
  \begin{tabularx}{\textwidth}{|Y|L{2.6cm}|}
    \hline
    \textbf{Funcionalidade (enunciado, Tema 6)} & \textbf{Requisitos} \\ \hline
    Pesquisa de produtos & RF01 \\ \hline
    Consulta a diferentes lojas independentes & RF02 \\ \hline
    Apresentação dos preços & RF05 \\ \hline
    Ordenação por preço & RF06 \\ \hline
    Identificação da loja com menor preço & RF07 \\ \hline
    Comparação entre consultas sequenciais e concorrentes & RF03, RF04, RF08 \\ \hline
    Evolução: comparação de características técnicas, avaliações e custo-benefício & RF05, RF06, RF07 \\ \hline
  \end{tabularx}
  \fonteautores
\end{quadro}

% -----------------------------------------------------------------
\chapter{Atores}

Os atores do sistema estão no \autoref{tab:atores}. O \textbf{agregador} é um componente interno do sistema, não um ator: ele intermedia o usuário e as lojas.

\begin{quadro}[!htbp]
  \centering\small
  \caption{Atores do sistema}\label{tab:atores}
  \begin{tabularx}{\textwidth}{|L{3.6cm}|L{3.4cm}|Y|}
    \hline
    \textbf{Ator} & \textbf{Tipo} & \textbf{Descrição} \\ \hline
    Usuário comprador & Primário (humano) & Informa o produto, o modo de consulta e o critério de ordenação, e analisa o resultado. \\ \hline
    Pesquisador / avaliador & Primário (humano) & Executa os experimentos de desempenho e de falha e analisa tabelas e gráficos. \\ \hline
    Serviço de loja & Secundário (sistema) & Processo independente que responde consultas sobre o seu catálogo. \\ \hline  \end{tabularx}
  \fonteautores
\end{quadro}

% -----------------------------------------------------------------
\chapter{Requisitos funcionais}

Os oito requisitos funcionais (\autoref{tab:rf}) seguem os exemplos da seção~5.2 do enunciado \cite{enunciado2026}, adaptados ao tema. As duas interfaces do usuário -- o \cod{cliente.py} de linha de comando e a interface web, que usa um gateway HTTP como cliente TCP do agregador -- atendem esses mesmos requisitos.

\legendaquadro{Requisitos funcionais}\label{tab:rf}\nopagebreak
\begin{xltabular}{\textwidth}{|L{1.2cm}|Y|L{1.9cm}|}
  \hline
  \textbf{ID} & \textbf{Requisito} & \textbf{Prioridade} \\ \hline
  \endfirsthead
  \multicolumn{3}{l}{\small\itshape (continuação)}\\
  \hline
  \textbf{ID} & \textbf{Requisito} & \textbf{Prioridade} \\ \hline
  \endhead
  \small RF01 & \small O sistema deverá permitir que o usuário informe os parâmetros de uma pesquisa: o produto (com a lista dos produtos disponíveis no catálogo das lojas), o modo de consulta (sequencial, concorrente ou comparação dos dois) e o critério de ordenação. & \small Essencial \\ \hline
  \small RF02 & \small O sistema deverá consultar diferentes serviços independentes: pelo menos quatro lojas, cada uma executando como processo servidor próprio. & \small Essencial \\ \hline
  \small RF03 & \small O sistema deverá permitir consultas \textbf{sequenciais}, em que uma loja é consultada após a resposta da anterior. & \small Essencial \\ \hline
  \small RF04 & \small O sistema deverá permitir consultas \textbf{concorrentes}, em que todas as lojas são consultadas simultaneamente por threads. & \small Essencial \\ \hline
  \small RF05 & \small O sistema deverá apresentar os resultados encontrados: para cada loja, o produto, as especificações técnicas, o preço, o frete, o preço total, a avaliação, a garantia e o tempo de resposta, além das lojas sem o produto e das que falharam, com o motivo (fora do ar, queda, timeout), sem deixar de apresentar as demais. & \small Essencial \\ \hline
  \small RF06 & \small O sistema deverá ordenar os resultados conforme um critério definido: preço total (preço + frete) ou índice de custo-benefício, à escolha do usuário. & \small Essencial \\ \hline
  \small RF07 & \small O sistema deverá identificar a melhor opção segundo o critério escolhido: a \textbf{loja com menor preço total} e a \textbf{oferta de melhor custo-benefício}. & \small Essencial \\ \hline
  \small RF08 & \small O sistema deverá registrar o tempo de execução das consultas (tempo total e tempo por loja) em \cod{resultados/log\_pesquisas.csv} (colunas \cod{tempo\_total\_ms} e \cod{tempos\_lojas\_ms}) e exibir o comparativo entre os modos sequencial e concorrente, com o speedup obtido. & \small Essencial \\ \hline\end{xltabular}
\fonteautores

% -----------------------------------------------------------------
\chapter{Requisitos não funcionais}

Os cinco requisitos não funcionais (\autoref{tab:rnf}) seguem os exemplos da seção~5.3 do enunciado \cite{enunciado2026}.

\legendaquadro{Requisitos não funcionais}\label{tab:rnf}\nopagebreak
\begin{xltabular}{\textwidth}{|L{1.4cm}|L{2.6cm}|Y|}
  \hline
  \textbf{ID} & \textbf{Categoria} & \textbf{Requisito} \\ \hline
  \endfirsthead
  \multicolumn{3}{l}{\small\itshape (continuação)}\\
  \hline
  \textbf{ID} & \textbf{Categoria} & \textbf{Requisito} \\ \hline
  \endhead
  \small RNF01 & \small Comunicação & \small A comunicação entre os componentes (cliente, agregador e lojas) deverá utilizar \textbf{sockets TCP} reais, com o protocolo de aplicação documentado (JSON UTF-8 terminado por \cod{\textbackslash n}, uma requisição por conexão, limite de 64~KB). \\ \hline
  \small RNF02 & \small Concorrência & \small A estratégia concorrente deverá utilizar \textbf{threads} (\cod{threading}), uma por loja consultada, de modo que o tempo total fique próximo da maior latência entre as lojas (sobrecarga inferior a 50~ms em ambiente local). \\ \hline
  \small RNF03 & \small Escalabilidade & \small O servidor deverá atender \textbf{múltiplas conexões simultaneamente}: o agregador e as lojas usam uma thread por conexão e devem atender pelo menos 50 clientes simultâneos sem falhas. \\ \hline
  \small RNF04 & \small Confiabilidade & \small O sistema deverá tratar falhas de comunicação (conexão recusada, queda durante a consulta, timeout de 2,5~s por loja, mensagens malformadas) sem encerrar o agregador e sem exibir \textit{traceback} ao usuário. \\ \hline
  \small RNF05 & \small Experimentação & \small A aplicação deverá permitir experimentos de desempenho: latência das lojas configurável (por loja ou por requisição), execução automatizada dos cinco experimentos com tabelas e gráficos e resultados reprodutíveis. \\ \hline\end{xltabular}
\fonteautores

% -----------------------------------------------------------------
\chapter{Restrições}

\begin{itemize}
  \item Linguagem \textbf{Python 3.10+} e apenas a biblioteca padrão (\cod{socket}, \cod{threading}, \cod{json}); \cod{matplotlib} só para os gráficos. Configuração centralizada em \cod{config.py}.
  \item Comunicação obrigatoriamente por \textbf{sockets TCP} \cite{kurose2021,python-socket}. Não é permitido simular a comunicação.
  \item Pelo menos \textbf{quatro} serviços independentes, cada um em processo separado e em porta TCP distinta (9001--9004). O agregador usa a porta 8000.
  \item Dados de produtos e preços \textbf{simulados}, armazenados em \cod{backend/dados/lojas/*.json}.
  \item Organização: o sistema distribuído (Python) fica em \cod{backend/} e a interface web em \cod{frontend/}. HTTP existe apenas entre o navegador e o gateway; o gateway é um cliente TCP do agregador, como o \cod{cliente.py}.
  \item Execução local (\cod{127.0.0.1}), com demonstração prática ao vivo na apresentação de 16/11/2026.
\end{itemize}

% -----------------------------------------------------------------
\chapter{Casos de uso}

A \autoref{fig:casos-uso} apresenta o diagrama UML de casos de uso; o documento de Modelagem descreve a notação e as relações \textit{include} e \textit{extend}. As seções a seguir detalham cada caso de uso.

\begin{figure}[htb]
  \centering
  \caption{Diagrama de casos de uso}\label{fig:casos-uso}
  \diagrama{casos_de_uso}
  \fonteautores
\end{figure}

\section{UC01 -- Pesquisar produto}

\begin{description}[leftmargin=0pt, labelsep=0.4em, itemsep=0.15em]
  \item[Ator:] Usuário comprador. \textbf{Requisitos:} RF01--RF08, RNF01--RNF04.
  \item[Pré-condição:] agregador e lojas em execução.
  \item[Fluxo principal:]\leavevmode
  \begin{enumerate}
    \item O usuário informa o produto, o modo e o critério.
    \item O cliente envia a requisição \cod{consultar} ao agregador.
    \item O agregador consulta as quatro lojas (sequencial ou concorrente).
    \item Cada loja responde com a sua oferta.
    \item O agregador ordena, identifica o menor preço e o melhor custo-benefício e registra o tempo.
    \item O cliente exibe o relatório.
  \end{enumerate}
  \item[Fluxos alternativos:]\leavevmode
  \begin{itemize}
    \item[A1:] produto inexistente $\rightarrow$ as lojas respondem \cod{nao\_encontrado} e o cliente informa que não há ofertas.
    \item[A2:] uma loja falha (fora do ar, queda ou timeout) $\rightarrow$ a falha é listada e as demais ofertas são exibidas.
    \item[A3:] agregador indisponível $\rightarrow$ o cliente exibe mensagem amigável.
    \item[A4:] parâmetro inválido $\rightarrow$ o agregador responde \cod{erro\_requisicao} e o cliente exibe o motivo.
  \end{itemize}
\end{description}

\section{UC02 -- Comparar estratégias sequencial × concorrente}

\begin{description}[leftmargin=0pt, labelsep=0.4em, itemsep=0.15em]
  \item[Ator:] Usuário comprador. \textbf{Requisitos:} RF03, RF04, RF08.
  \item[Fluxo:] o usuário escolhe o modo ``comparar''. O cliente executa o UC01 nos dois modos e exibe o relatório de cada um, os dois tempos e o speedup.
\end{description}

\section{UC03 -- Listar catálogo}

\begin{description}[leftmargin=0pt, labelsep=0.4em, itemsep=0.15em]
  \item[Ator:] Usuário comprador. \textbf{Requisito:} RF01.
  \item[Fluxo:] o cliente envia \cod{listar\_produtos}. O agregador consulta as lojas concorrentemente e devolve a união das palavras-chave.
\end{description}

\section{UC04 -- Executar experimentos}

\begin{description}[leftmargin=0pt, labelsep=0.4em, itemsep=0.15em]
  \item[Ator:] Pesquisador. \textbf{Requisitos:} RF08, RNF03, RNF05.
  \item[Fluxo:] \cod{experimentos.py} sobe as lojas e o agregador como processos, executa os experimentos 1 a 5 com repetições e grava os CSVs e \cod{tabelas.md}. Em seguida, \cod{graficos.py} gera os gráficos.
  \item[Fluxo alternativo (ao vivo, na apresentação):] com \cod{iniciar\_servicos.py} no ar, o pesquisador executa cada experimento pelo \cod{cliente.py}: \cod{--modo sequencial} (Exp.~1), \cod{--modo concorrente} (Exp.~2), \cod{--atraso MS} (Exp.~3), \cod{--clientes N} (Exp.~4) e, para o Exp.~5, digita \cod{derrubar <loja>} no terminal do \cod{iniciar\_servicos.py} durante uma pesquisa e depois \cod{subir <loja>}.
\end{description}

\section{UC05 -- Usar a interface web}

\begin{description}[leftmargin=0pt, labelsep=0.4em, itemsep=0.15em]
  \item[Ator:] Usuário comprador ou avaliador. \textbf{Requisitos:} RF01--RF08, RNF03--RNF05 (a interface é mais um cliente do agregador).
  \item[Pré-condição:] \cod{backend/gateway\_http.py} em execução (ele sobe as lojas e o agregador) e a interface aberta no navegador.
  \item[Fluxo principal:] o usuário pesquisa na aba \textit{Pesquisa}; o gateway envia \cod{consultar} ao agregador por TCP (duas vezes no modo ``comparar'') e a interface mostra destaques, tempos por loja, speedup e a tabela de ofertas.
  \item[Fluxos alternativos:]\leavevmode
  \begin{itemize}
    \item[A1:] na aba \textit{Lojas e falhas}, o usuário derruba uma loja; a pesquisa seguinte mostra 3/4 lojas e o motivo da falha; ao subir a loja de novo, ela volta a responder.
    \item[A2:] na aba \textit{Teste de carga}, o usuário dispara N clientes TCP simultâneos (1--50) e vê média, p95, vazão e falhas.
    \item[A3:] gateway ou agregador fora do ar $\rightarrow$ a interface exibe mensagem amigável, sem \textit{traceback}.
  \end{itemize}
\end{description}

% -----------------------------------------------------------------
\chapter{Critérios de aceitação}

Cada critério do \autoref{tab:ca} pode ser verificado por um experimento ou por um teste automatizado da pasta \cod{backend/tests/}. A suíte tem 90 testes e todos passam.

{\small
\legendaquadro{Critérios de aceitação}\label{tab:ca}\nopagebreak
\begin{xltabular}{\textwidth}{|L{1.0cm}|Y|L{5.6cm}|}
  \hline
  \textbf{ID} & \textbf{Critério} & \textbf{Verificação} \\ \hline
  \endfirsthead
  \multicolumn{3}{l}{\itshape (continuação)}\\
  \hline
  \textbf{ID} & \textbf{Critério} & \textbf{Verificação} \\ \hline
  \endhead
  CA01 & Uma pesquisa por ``rtx 4060'' com as 4 lojas ativas retorna 4 ofertas ordenadas por preço total crescente, e o menor preço é da PichauSim (R\$~1.864,00). & \cod{test\_ordena\_por\_preco\_e\_identifica\_menor}, \cod{test\_pesquisa\_via\_cliente} \\ \hline
  CA02 & Com o critério \cod{custo\_beneficio}, as ofertas aparecem em ordem decrescente de score. & \cod{test\_ordena\_por\_custo\_beneficio} \\ \hline
  CA03 & Os modos sequencial e concorrente produzem o mesmo resultado (mesma ordem de lojas). & \cod{test\_sequencial\_e\_concorrente\_dao\_mesmo\_resultado} \\ \hline
  CA04 & Com 4 lojas de 100~ms, o modo sequencial leva $\geq$~400~ms e o concorrente $<$~250~ms (speedup $>$~2,5). & \cod{test\_concorrente\_mais\_rapido\_que\_sequencial}, Exp.~1--3 \\ \hline
  CA05 & Com lojas heterogêneas (50/100/150/200~ms), o modo sequencial $\approx$ soma e o concorrente $\approx$ máximo das latências. & \cod{test\_latencia\_heterogenea}, Exp.~3b \\ \hline
  CA06 & 20 clientes simultâneos recebem respostas completas, sem falhas. Nos experimentos, 50 clientes. & \cod{test\_multiplos\_clientes\_simultaneos}, Exp.~4 \\ \hline
  CA07 & Com uma loja desligada, a pesquisa retorna 3 ofertas e aponta \cod{erro\_comunicacao} para a loja ausente. & \cod{test\_loja\_desligada}, Exp.~5a \\ \hline
  CA08 & Se a loja cair durante a consulta, a falha é detectada imediatamente (sem esperar o timeout) e a pesquisa conclui com 3 ofertas. & \cod{test\_loja\_cai\_durante\_a\_consulta}, Exp.~5b \\ \hline
  CA09 & Uma loja que não responde em 2,5~s recebe status \cod{timeout}, e a pesquisa concorrente não passa de $\approx$~2,5~s. & \cod{test\_loja\_lenta\_gera\_timeout}, Exp.~5c \\ \hline
  CA10 & Requisições inválidas (operação desconhecida, produto vazio, modo/critério/atraso inválidos, JSON malformado) recebem \cod{erro\_requisicao}, e o servidor continua atendendo. & \cod{test\_requisicoes\_invalidas}, \cod{test\_json\_malformado\_recebe\_erro}, \cod{test\_erro\_interno\_nao\_derruba\_servidor} \\ \hline
  CA11 & Cada pesquisa gera uma linha em \cod{log\_pesquisas.csv} com o tempo total e o tempo de cada loja. & \cod{test\_registra\_tempo\_em\_csv}, \cod{test\_registra\_tempo\_de\_cada\_loja} \\ \hline
  CA12 & Com o agregador desligado, o cliente exibe uma mensagem de erro amigável e termina com código 1. & \cod{test\_agregador\_fora\_do\_ar} \\ \hline
  CA13 & Uma loja que responde \cod{sucesso} com dados incompletos ou de tipo errado é registrada como \cod{erro\_comunicacao}, e a pesquisa conclui com as demais lojas. & \cod{test\_loja\_com\_resposta\_incompleta}, \cod{test\_loja\_com\_preco\_nao\_numerico} \\ \hline
  CA14 & Com as portas já ocupadas, \cod{iniciar\_servicos.py} recusa a subida (código 1) em vez de anunciar um ambiente que não é o seu; ao receber Ctrl+C, SIGTERM ou SIGHUP, encerra todas as lojas e o agregador. & \cod{test\_porta\_ocupada\_e\_detectada}, \cod{test\_6\_segundo\_ambiente\_com\_portas\_ocupadas\_falha}, \cod{test\_sigterm\_encerra\_todos\_os\_processos} \\ \hline
  CA15 & O roteiro da demonstração segue os 5 experimentos do enunciado e funciona com processos reais nas portas 8000 e 9001--9004: sequencial, concorrente, atraso artificial, N clientes simultâneos (\cod{cliente.py --clientes N}) e uma loja derrubada durante a pesquisa pelo console do \cod{iniciar\_servicos.py} (\cod{derrubar}/\cod{subir}), além de listar, custo-benefício e parâmetro inválido. & \cod{tests/test\_sistema.py} (\cod{TestDemonstracao}), \cod{test\_console\_derruba\_e\_sobe\_loja}, \cod{test\_cli\_multiplos\_clientes} \\ \hline
  CA16 & Pelo gateway HTTP, a pesquisa ``rtx 4060'' retorna o menor preço da PichauSim (R\$~1.864,00), o modo ``comparar'' devolve os dois modos e o speedup, e requisições inválidas recebem HTTP 400 \cod{erro\_requisicao}. & \cod{test\_gateway.py} (\cod{test\_consulta\_concorrente}, \cod{test\_comparar\_devolve\_os\_dois\_modos\_e\_speedup}, \cod{test\_requisicoes\_invalidas}) \\ \hline
  CA17 & Pelo gateway, uma loja pode ser derrubada e subida de novo (processo real); com uma loja fora do ar a pesquisa retorna 3/4 lojas; com o agregador fora do ar a resposta é HTTP 502 com mensagem amigável. & \cod{test\_derrubar\_e\_subir\_loja}, \cod{test\_loja\_desligada\_aparece\_na\_consulta}, \cod{test\_agregador\_fora\_do\_ar} \\ \hline
  CA18 & O teste de carga do gateway atende 20 clientes simultâneos sem falhas e recusa mais de 50; o gateway só serve gráficos de \cod{resultados/} (sem acesso a outros arquivos). & \cod{test\_carga\_com\_20\_clientes}, \cod{test\_carga\_fora\_do\_limite}, \cod{test\_resultados\_nao\_permite\_sair\_do\_diretorio} \\ \hline\end{xltabular}}
\fonteautores

% -----------------------------------------------------------------
\chapter{Rastreabilidade}

O \autoref{tab:rastreabilidade} liga cada requisito aos casos de uso e aos módulos que o implementam. Os módulos Python ficam em \cod{backend/}; a interface web, em \cod{frontend/}.

{\small
\legendaquadro{Matriz de rastreabilidade}\label{tab:rastreabilidade}\nopagebreak
\begin{xltabular}{\textwidth}{|L{1.9cm}|L{3.4cm}|Y|}
  \hline
  \textbf{Requisito} & \textbf{Casos de uso} & \textbf{Módulos} \\ \hline
  \endfirsthead
  \multicolumn{3}{l}{\itshape (continuação)}\\
  \hline
  \textbf{Requisito} & \textbf{Casos de uso} & \textbf{Módulos} \\ \hline
  \endhead
  RF01 & UC01, UC02, UC03, UC05 & \cod{cliente.py}, \cod{agregador.py} (\cod{listar\_produtos}), \cod{gateway\_http.py} + \cod{frontend/src/} \\ \hline
  RF02 & UC01 & \cod{loja\_server.py}, \cod{config.py}, \cod{iniciar\_servicos.py} \\ \hline
  RF03, RF04 & UC01, UC02 & \cod{agregador.py} (\cod{consultar\_todas}) \\ \hline
  RF05 & UC01, UC05 & \cod{agregador.py} (\cod{executar\_busca}), \cod{cliente.py}, \cod{frontend/src/} \\ \hline
  RF06, RF07 & UC01 & \cod{agregador.py} (\cod{executar\_busca}, \cod{calcular\_score\_custo\_beneficio}) \\ \hline
  RF08 & UC01, UC02, UC04 & \cod{agregador.py} (\cod{RegistroPesquisas}), \cod{cliente.py} (\cod{comparar}) \\ \hline
  RNF01 & todos & \cod{protocolo.py}, \cod{servidor\_base.py} \\ \hline
  RNF02 & UC01, UC02 & \cod{agregador.py} (uma thread por loja + \cod{join()}) \\ \hline
  RNF03 & UC01, UC04 & \cod{servidor\_base.py}, \cod{cliente.py} (\cod{disparar\_clientes\_simultaneos}, \cod{--clientes}) \\ \hline
  RNF04 & UC01 & \cod{protocolo.py}, \cod{agregador.py} (\cod{validar\_resposta\_loja}), \cod{cliente.py}, \cod{iniciar\_servicos.py} (\cod{derrubar}/\cod{subir}) \\ \hline
  RNF05 & UC04 & \cod{loja\_server.py} (latência), \cod{protocolo.py} (\cod{atraso\_ms}), \cod{iniciar\_servicos.py} (\cod{--latencias}), \cod{experimentos.py}, \cod{graficos.py} \\ \hline\end{xltabular}}
\fonteautores

% ================================================================= PÓS-TEXTUAIS
\postextual
\bibliography{referencias}

\end{document}
